%% Revised manuscript (PeerJ CS-2026:06:142937, major revision).
%% Every number is generated: \R{key} from generated/numbers*.tex, tables via \input. Missing keys
%% render as red [key]. Sources are the frozen outputs in revision/tables/families/.
\documentclass[fleqn,10pt,lineno]{wlpeerj}
%% T1 is required under XeTeX (tectonic): with the default TU encoding the template's Times falls
%% back to Latin Modern regular, which loses bold and italics and can crash xdvipdfmx.
\usepackage[T1]{fontenc}
\usepackage{bm}
\usepackage{multirow}
\usepackage{url}
\usepackage{booktabs}
\usepackage{xcolor}
\usepackage{amsmath}
\usepackage{graphicx}

\makeatletter
\input{generated/numbers.tex}
\input{generated/numbers_extra.tex}
\input{generated/numbers_expansion.tex}
\newcommand{\R}[1]{\ifcsname R@#1\endcsname\csname R@#1\endcsname\else\textbf{\textcolor{red}{[#1]}}\fi}
\makeatother
\newcommand{\sg}{\operatorname{sg}}

\title{When does boundary distillation help lightweight medical image segmentation?}

\author[1]{Jingchen Chen}
\author[1]{Chenxiang Hong}
\author[1]{Yao Wu}
\author[1]{Tao Yang}
\affil[1]{Future Technology School, Shenzhen Technology University, Shenzhen, Guangdong, China}
\corrauthor[1]{Tao Yang}{yangtao@sztu.edu.cn}

\begin{abstract}
A boundary decoder attached during training and discarded afterwards gives a compact segmentation network contour supervision at no inference cost. Boundary distillation goes further and tries to move what that decoder learns into the deployed network. We ask whether this step adds anything beyond the simple supervision. BDS-Lite is a recipe for a 2D U-Net in which a removable sibling decoder teaches the segmentation features through stop-gradient distillation and a bounded gate, combined with a signed-distance surface loss, so that each signal can be removed in turn. We compared it with six ablations and four comparators under one protocol on ISIC2018, ACDC, and Synapse. Because the expected effects are as small as retraining variation, recipes were compared across independent runs, and ISIC2018 contrasts were confirmed on fresh seeds. An audit showed that the ISIC2018 results of our earlier submission were measured largely on training images; they are withdrawn and replaced by the official 1,000-image test set. With ten runs per arm on ISIC2018, the decoder alone scored above U-Net in Boundary F1 (\R{X.decoder.unet.boundary_f1.all.delta}, 95\% interval \R{X.decoder.unet.boundary_f1.all.ci}; supported, not confirmed on fresh seeds) with no relevant change in DSC or the 95th-percentile Hausdorff distance (HD95). BDS-Lite (\R{X.full.decoder.boundary_f1.all.delta}, \R{X.full.decoder.boundary_f1.all.ci}) and U-Net with Generalized Surface Loss (\R{X.gsl.decoder.boundary_f1.all.delta}, \R{X.gsl.decoder.boundary_f1.all.ci}) both scored below the decoder alone, and both differences were confirmed. Adding distillation and the gate to the decoder gave no established benefit: pooled over datasets with three runs per arm, Boundary F1 changed by \R{I.distillationandgate-without-surface.pooled.boundary_f1.three.est} \R{I.distillationandgate-without-surface.pooled.boundary_f1.three.ci} without and \R{I.distillationandgate-with-surface.pooled.boundary_f1.three.est} \R{I.distillationandgate-with-surface.pooled.boundary_f1.three.ci} with the surface loss, and DSC and HD95 showed no relevant effect. Distillation placed boundary signal in a deployed projection (average precision \R{MECH.isic2018.head.dist} versus \R{MECH.isic2018.head.plain} on ISIC2018) but not detectably in the features that make the prediction. Identical recipes ``differed'' at $p<0.05$ in about \R{SS.same.mean.3.boundary_f1.p05}\% of image-level comparisons, and three-seed results that had passed multiple-testing correction did not survive further runs. In the studied setting, boundary distillation showed no established value beyond training-only boundary supervision, and effects of this size can be judged only over independent runs.
\end{abstract}

\begin{document}
\flushbottom
\maketitle
\thispagestyle{empty}

%% =====================================================================
\section*{Introduction}

Shifting a contour by one pixel changes the Dice similarity coefficient (DSC) by about \R{S.isic2018.1.drop} on an average ISIC2018 skin lesion and by about \R{S.acdc.2.drop} on the ACDC myocardium (Table~\ref{tab:datasets}). Overlap scores therefore say little about contour placement on large objects, even though the contour is what several downstream measurements use. Border irregularity is one criterion of the ABCD rule for melanoma screening \citep{Nachbar1994ABCD}; ejection fraction and myocardial mass are computed from delineated cardiac volumes \citep{Bernard2018Deep}; organ contours define volumes and organs at risk in abdominal radiotherapy planning \citep{Jabbour2014Upper}. Boundary F1, the 95th-percentile Hausdorff distance (HD95) \citep{Huttenlocher1993Comparing}, and the average symmetric surface distance (ASSD) measure contour error directly, and evaluation guidance recommends reporting them alongside overlap \citep{Taha2015Metrics, Maier-Hein2024Metrics}.

Several families of methods target these errors during training. Distance-based losses add a geometric penalty to the output probabilities \citep{Kervadec2019Boundary, Karimi2020Reducing, Celaya2024A, Sun2023Boundary, Wang2022Active}. Auxiliary branches predict contours or distance maps from shared features \citep{Murugesan2019PsiNet, Takikawa2019GatedSCNN}; the branch is often kept at inference, as in the lightweight LB-UNet \citep{Xu2024LBUNet}, but it can also be discarded, as STDC does with a detail head trained on boundary-derived targets \citep{Fan2021STDC}. This last option is attractive for compact models, which remain in active use \citep{Ruan2023EGEUNet, Jiang2024LVUNet}, because boundary supervision then costs nothing at deployment. Boundary distillation goes one step further. It treats the boundary branch as a teacher and tries to move its knowledge into the network that is kept \citep{Liu2024BPKD, Wang2023MISSU, Banerjee2024Dual}.

That extra step raises a question that our title asks loosely and the literature seldom asks at all: help relative to what? A boundary branch that shares an encoder with the segmentation path already shapes the deployed features, so the proper reference for boundary distillation is the same network trained with the same branch and no transfer. Reports usually compare against a network with no boundary signal at all, and composite recipes are rarely decomposed, so this contrast is seldom available. A second gap is measurement. If the effects are similar in size to the variation between retrainings of one recipe, a single training per arm, or a test over the images of one trained model, cannot show them.

This paper addresses both gaps. BDS-Lite adds to a compact U-Net \citep{Ronneberger2015UNet} a boundary decoder that is removed after training. Its features serve as a stop-gradient distillation target for the sibling segmentation decoder, a small bounded gate lets the distilled features rescale the output, and a signed-distance loss targets distant errors. What distinguishes the recipe is a specific combination: the teacher is a sibling decoder of one compact 2D network, not a separate model, and every signal can be switched off, so that the decoder-only arm is exactly the simple supervision against which distillation has to be judged.

We hold the backbone and training protocol fixed. The compact setting is deliberate: removing components at deployment matters most under tight budgets, a fixed small backbone keeps the comparison controlled, and every experiment fits on one 8\,GB laptop GPU. The tasks differ in structure: one lesion per image (ISIC2018), three adjacent cardiac structures (ACDC), and eight abdominal organs of very different size (Synapse). Measurement is part of the design: we analyze recipes over independent runs and tested the ISIC2018 contrasts on seeds that had not suggested them. Tracing our earlier results to their checkpoints also showed that the ISIC2018 models of our original submission had been evaluated largely on their own training images, and this paper corrects that evaluation.

The evidence points to a clear ordering. Simple training-only supervision gave the best-supported favorable result: among the four ISIC2018 arms trained ten times it had the highest Boundary F1, a gain over U-Net that is supported but not confirmed, with the same direction on the other datasets and no cost at inference. Distillation and the gate added no established benefit to it: distillation placed boundary-readable signal in a deployed projection, but that signal did not detectably reach the features that make the prediction. The full recipe scored below the decoder alone on ISIC2018, a confirmed result. Our contributions are:
\begin{enumerate}
\item a direct test of whether boundary distillation adds to training-only boundary supervision, using a matched ablation with a decoder-only reference and four comparators on three tasks;
\item a modular recipe, BDS-Lite, in which each boundary signal can be removed and whose deployed graph is specified exactly, with its deployment and training costs;
\item direct measurements of what distillation and the gate do inside the network, which show how a transfer can succeed by its own objective and leave the prediction unchanged;
\item a recipe-level evaluation for effects of this size, with run-level intervals, confirmation on fresh seeds, a quantification of how retraining variation misleads image-level tests, and a corrected ISIC2018 evaluation.
\end{enumerate}

%% =====================================================================
\section*{Related work}

\emph{Boundary losses.} Boundary loss integrates softmax probabilities against a signed distance map and needs a schedule that rebalances it with a region loss \citep{Kervadec2019Boundary}. Hausdorff-oriented losses penalize large surface deviations at the cost of heavier distance computation \citep{Karimi2020Reducing}. Generalized Surface Loss (GSL) normalizes a distance-weighted error and anneals from a region term to a surface term \citep{Celaya2024A}; Boundary DoU measures mismatch inside a boundary band \citep{Sun2023Boundary}; active boundary loss aligns predicted and reference boundaries during training \citep{Wang2022Active}. None of these adds inference cost, but their reported gains come from different backbones and protocols, so their relative merit is hard to judge from the literature.

\emph{Boundary branches and compact networks.} Psi-Net predicts mask, contour, and distance map jointly \citep{Murugesan2019PsiNet}, and Gated-SCNN routes shape information through a gated stream \citep{Takikawa2019GatedSCNN}; such pathways are often kept at inference or tied to backbone changes. Lightweight networks reduce cost by other means, such as grouped attention in EGE-UNet \citep{Ruan2023EGEUNet} and MLP blocks in UNeXt \citep{Valanarasu2022UNeXt}, usually with overlap-only evaluation. LB-UNet uses boundary prediction modules, but they stay in the deployed model and the network was designed for binary lesions \citep{Xu2024LBUNet}. Discarding a boundary head after training is not new: STDC does so with a detail head trained with binary cross-entropy and Dice \citep{Fan2021STDC}, which is the pattern of our decoder-only arm. Heavier systems such as nnU-Net \citep{Isensee2021nnUNet}, TransUNet \citep{Chen2021TransUNet}, and Swin-Unet \citep{Cao2022SwinUnet} gain accuracy from capacity, pretraining, or extensive augmentation, and lie outside the compact regime studied here.

\emph{Distillation.} Knowledge distillation trains a student on signals produced by a teacher \citep{Hinton2015Distilling}. Feature-based variants match intermediate representations, and hint training places a learned regressor between student and teacher because their shapes differ \citep{Romero2015FitNets}; this detail matters for our results. In segmentation, distillation has transferred structured outputs \citep{Liu2019Structured} and boundary and body knowledge from a separate teacher \citep{Liu2024BPKD}, and self-distillation within one network has been applied to 3D medical segmentation \citep{Wang2023MISSU, Banerjee2024Dual}. Such evaluations leave open how much comes from the transfer itself and how much from the extra supervision. In BDS-Lite the teacher is a sibling decoder of the same compact 2D network, detached and discarded, and the design isolates this contrast.

\emph{Measuring small effects.} Multitask objectives can interfere in a shared encoder \citep{Yu2020PCGrad}, so an auxiliary task need not help its main task, and a recipe with several parts needs ablation before any part is credited. For comparing several methods, rank-based omnibus tests with corrected post-hoc tests are recommended \citep{Demsar2006Statistical}. Such tests treat test samples as the unit, which suits fixed models; our data address whether it suits training recipes. Supplementary S1 tabulates mechanism, inference cost, reported metrics, and limitations of the methods above.

%% =====================================================================
\section*{Materials \& Methods}

\subsection*{BDS-Lite}
The backbone is a four-level 2D U-Net with base width $c=32$ (1.93\,M parameters). The encoder produces multi-scale features, the segmentation decoder $D_s$ produces a final feature map $F_s\in\mathbb{R}^{c\times H\times W}$, and a $1\times1$ convolution $W_o$ gives logits for $K$ classes. During training, an auxiliary boundary decoder $D_b$ with three upsampling blocks reads the same encoder features and outputs a boundary feature $q_b\in\mathbb{R}^{c_b\times H\times W}$ ($c_b=16$) and a single-channel boundary logit map $z_b$ (Fig.~\ref{fig:architecture}). A $1\times1$ convolution $P$ maps the segmentation feature to $q_s=P(F_s)$ with the shape of $q_b$, and $q_s$ drives a gate:
\begin{equation}
g=\sigma\!\left(W_g q_s+b_g\right)\in(0,1)^{c\times H\times W},\qquad
F_s^{\mathrm{m}}=F_s\odot\bigl[1+\alpha\,(g-0.5)\bigr],\qquad z_s=W_oF_s^{\mathrm{m}},
\label{eq:gate}
\end{equation}
where $W_g\in\mathbb{R}^{c\times c_b}$ is a $1\times1$ convolution, so the gate is per-pixel and per-channel. Both $W_g$ and $b_g$ start at zero, so training begins from the identity. Because $g\in(0,1)$, the factor lies in $[1-\alpha/2,\,1+\alpha/2]$: $\alpha$ bounds how much the gate can rescale a feature, and $\alpha=0$ removes it. We fixed $\alpha=0.25$ ($\pm12.5\%$) before the experiments to keep the network near the plain U-Net; it was not tuned on test data. At deployment $D_b$ and all auxiliary losses are deleted, which leaves the output unchanged because $q_b$ and $z_b$ do not feed the segmentation path. The projection $P$ and the gate remain ($c\,c_b+c_b+c_b\,c+c=1{,}072$ parameters).

\begin{figure}[ht]
\centering
\includegraphics[width=\linewidth]{figure_architecture_v2}
\caption{BDS-Lite. Blue: deployed components (encoder, segmentation decoder $D_s$, projection, gate, output convolution). Orange: training-only boundary decoder $D_b$ and losses. The boundary target is the morphological gradient (dilation XOR erosion) of the class masks. $L_{\mathrm{bnd}}$ and $L_{\mathrm{sdf\text{-}surface}}$ are $L_{\mathrm{boundary}}$ and $L_{\mathrm{sdf}}$ in the text, and $\lambda_{\mathrm{seg}}=1$. The decoder-only arm keeps the orange decoder with $L_{\mathrm{bnd}}$ and has no $L_{\mathrm{distill}}$ or gate; its deployed network is a plain U-Net.}
\label{fig:architecture}
\end{figure}

Let $p=\mathrm{softmax}(z_s)$ and let $y$ be the one-hot reference. The segmentation loss is $L_{\mathrm{seg}}=L_{\mathrm{ce}}+L_{\mathrm{dice}}$, with pixel-averaged cross-entropy and soft Dice computed per foreground class over the batch and averaged over classes. Each class thus contributes its own ratio regardless of its area, which counteracts background dominance for small structures; whether this suffices is examined in the Results.

The boundary target $y_b$ is the union over foreground classes of the morphological gradient, the exclusive-or of the binary dilation and erosion with a four-connected $3\times3$ element, and is two pixels thick. The decoder is trained with $L_{\mathrm{boundary}}=\mathrm{BCE}(z_b,y_b)+L_{\mathrm{Dice}}(\sigma(z_b),y_b)$. Distillation uses
\begin{equation}
L_{\mathrm{distill}}=\mathrm{SmoothL1}_{\beta=1}\bigl(q_s,\sg(q_b)\bigr),
\end{equation}
where the stop-gradient operator $\sg(\cdot)$ is the identity in the forward pass and has zero Jacobian in the backward pass (\texttt{detach()} in PyTorch). $L_{\mathrm{distill}}$ therefore updates $P$ and, through it, the segmentation decoder and encoder, but sends no gradient into $D_b$, which is trained only by $L_{\mathrm{boundary}}$. At each step the boundary feature is a fixed target, so the transfer runs one way.

For image $n$ and foreground class $k$, let $\phi_{nk}(u)$ be the signed Euclidean distance of pixel $u$ to the class boundary (positive outside; $\phi\equiv0$ if the class is absent), $s_{nk}=\max_u|\phi_{nk}(u)|$, and $v_{nk}=\mathbb{1}[s_{nk}>0]$. The surface loss is
\begin{equation}
L_{\mathrm{sdf}}=\frac{\sum_{n,k}v_{nk}\sum_u p_{nk}(u)\,\phi_{nk}(u)/s_{nk}}{HW\sum_{n,k}v_{nk}}.
\label{eq:sdf}
\end{equation}
Probability mass outside the object is penalized in proportion to its normalized distance from the boundary, and mass inside is rewarded, so the loss can be negative. The full objective is
\begin{equation}
L=L_{\mathrm{seg}}+\lambda_{\mathrm{sdf}}L_{\mathrm{sdf}}+\lambda_bL_{\mathrm{boundary}}+\lambda_dL_{\mathrm{distill}},\qquad \lambda_{\mathrm{sdf}}=1,\ \lambda_b=\lambda_d=0.05.
\label{eq:bds}
\end{equation}
The weights were fixed before the experiments and are identical across datasets; in the first epoch of the seed-1 diagnostic runs, the weighted boundary and distillation terms equaled \R{AUX.isic2018.share}\%, \R{AUX.acdc.share}\%, and \R{AUX.synapse.share}\% of $L_{\mathrm{seg}}$ on ISIC2018, ACDC, and Synapse. Supplementary S2 gives the training step as an algorithm.

\subsection*{How boundary information can reach the prediction}
Boundary supervision can affect the deployed output by three routes, and our question is whether the second and third add to the first.

\emph{Shared encoder.} $L_{\mathrm{boundary}}$ trains $D_b$, and its gradient continues into the encoder, whose features $D_s$ also reads. It reaches neither $D_s$ nor $W_o$. This is the only route available to the decoder-only arm, and it should act mainly on contour placement, hence on Boundary F1, because the loss concentrates on a two-pixel band.

\emph{Distillation.} $L_{\mathrm{distill}}$ is the only signal from the boundary decoder that reaches $D_s$. This gives the reason to expect a benefit: if the encoder holds boundary information that $D_s$ does not use, a target applied to the decoder's output feature could make use of it. It also gives a reason for doubt. As in hint training \citep{Romero2015FitNets}, student and teacher are joined by a learned map, here $P$, and the loss can fall because $P$ reads what $F_s$ already contains, with $F_s$ itself unchanged. The transfer would then succeed by its own objective without changing the prediction. The two cases can be told apart only by measuring $q_s$ and $F_s$ separately, which we do.

\emph{Gate.} The gate gives $q_s$ a direct path to the output. If it carried boundary information, its modulation should concentrate near contours; because it is bounded, removing it should change predictions only slightly.

The surface loss is a separate signal, unrelated to the transfer. It weights probability mass by distance from the reference contour and should reduce distant false positives, hence HD95 and ASSD. We wrote these expectations down with the earlier results in hand, so they explain the design and are tested by the ablations, but they are not preregistered hypotheses.

A first-order argument links contour error to the metrics. If a contour of length $L$ around an object of area $A$ is displaced by a mean normal offset $\delta$, the symmetric difference between prediction $P$ and reference $G$ is $|P\triangle G|\approx L|\delta|$, so
\begin{equation}
\mathrm{DSC}\approx1-\frac{L\,|\delta|}{2A},
\label{eq:firstorder}
\end{equation}
whereas ASSD $\approx|\delta|$ and HD95 follows the tail of the distance distribution. Table~\ref{tab:datasets} lists $L/(2A)$ for each dataset. On large lesions, contour placement can thus change with little effect on DSC, while distance metrics also respond to errors far from the contour that a band-limited boundary signal does not address. A boundary signal should therefore be expected to move Boundary F1 first, and this is the reason it is the primary endpoint.

\subsection*{Datasets}
ISIC2018 Task~1 \citep{Codella2019Skin} provides dermoscopic images with one lesion class. In \R{CMP.isic2018.single}\% of training masks the lesion is a single component. ACDC \citep{Bernard2018Deep} provides cine MRI with right ventricle, myocardium, and left ventricle. Synapse \citep{Landman2015MICCAI} provides abdominal CT with eight organs; its 12 test cases are the standard TransUNet test cases \citep{Chen2021TransUNetArxiv}, but we train on 14 rather than 18 cases because four are held out for validation. We chose these tasks because they differ in the number and size of structures and in how often a class is absent (Table~\ref{tab:datasets}, Fig.~\ref{fig:datasets}), and because contours serve a different purpose in each. They do not form a controlled scale of difficulty: they also differ in modality and sample size. We did not include LiTS or BraTS because they are 3D and, for BraTS, multimodal, which lies outside a compact single-channel 2D scope and our single-GPU budget; slice-wise transformer pipelines have been proposed for such data \citep{Raj2026ViTSwinUNet}. All inputs are $224\times224$ arrays produced by the released converters, with continuous interpolation for images and nearest-neighbor interpolation for masks. ISIC2018 images are scaled to $[0,1]$, and the provided ACDC and Synapse intensities are kept. The datasets are public, de-identified benchmarks used under their providers' terms; no new human data were collected and institutional review board approval was not required.

\begin{table}[ht]
\centering
\resizebox{\linewidth}{!}{\input{generated/tab_datasets.tex}}
\caption{\label{tab:datasets}Datasets. Class absent: fraction of training images without the class (range over classes). DSC loss per 1\,px shift: $L/(2A)$ from Eq.~\ref{eq:firstorder}, computed on non-empty training masks (range over classes). Per-class descriptors are in Supplementary S3.}
\end{table}

\begin{figure}[ht]
\centering
\includegraphics[width=0.78\linewidth]{figure_dataset_examples}
\caption{\label{fig:datasets}Training examples and derived targets: image, reference mask, union boundary target $y_b$, and signed-distance target $\phi$ (minimum over classes; negative inside). One sample per dataset was selected by a fixed rule (most classes present, then median foreground area). ISIC2018 image: ISIC archive, CC-0. ACDC and Synapse intensity images are not shown because we could not confirm that their terms of use permit republication.}
\end{figure}

\subsection*{Correction of the ISIC2018 evaluation}
The ISIC2018 results of our original submission are invalid and are withdrawn. That submission stated that U-Net and BDS-Lite were trained on 2,075 images and evaluated on 519 held-out images of the 2,594-image development pool. Tracing each number to its checkpoint showed otherwise. All six ISIC2018 checkpoints record 226.9 optimizer steps per epoch, which matches an earlier partition of 1,815 training and 389 validation images (batch size 8) and not the stated one, and their selection logs match the 389-image validation set. Of the 519 images used for the reported evaluation, 358 belonged to these models' training set, 84 to their validation set, and 77 to neither. The reported ISIC2018 values were therefore not held-out estimates, and the claimed HD95 benefit rested on them. ACDC and Synapse models used the documented partitions, and their original results reproduce with the present evaluator (Supplementary S4).

We replace the invalid evaluation with the official ISIC2018 test set (1,000 images with released reference masks, CC-0; \url{https://challenge.isic-archive.com/data/}), which shares no image identifier with the development pool and was never used for any selection. The ISIC archive does contain repeated images of some lesions under different identifiers: content matching found \R{DUP.n.0.99} test images with a near-identical development image ($r>0.99$) and up to \R{DUP.n.0.9} with $r>0.90$. Because no threshold separates repeats cleanly, we also repeat the main contrasts after excluding test images at four thresholds (Supplementary S4). Patient identities are unavailable, so this is an independent image set, not a patient-independent cohort.

\subsection*{Model sets and compared arms}
All runs used AdamW (weight decay $10^{-5}$), cosine annealing over 150 epochs, gradient-norm clipping at 1.0, automatic mixed precision, and no online augmentation, with batch size 8 and learning rate $3\times10^{-4}$ on ISIC2018 and batch size 16 and learning rate $10^{-4}$ on ACDC and Synapse. The checkpoint with maximum validation DSC was retained; test data were never used for selecting models, checkpoints, hyperparameters, or thresholds. Experiments ran on one NVIDIA RTX 5060 Laptop GPU (8\,GB) with CUDA 13.0 and PyTorch 2.11.0.

The \emph{original} set comprises the U-Net and BDS-Lite checkpoints of the first submission (seeds 1--3); on ISIC2018 they used the 1,815-image partition, and the pre-release trainer that produced them is not preserved. The \emph{retrained} set uses the released code on the documented partitions throughout (ISIC2018: 2,075 training, 519 validation images), with seeds 1--3 for every arm. On ISIC2018 its U-Net, BDS-Lite, surface-only, and BDS-Lite-without-surface runs are reused from a follow-up study on the same pipeline, whose runner reproduced the main U-Net and BDS-Lite models exactly (Supplementary S5). Hyperparameters are the same in both sets, and no arm in the analyses below uses class weights (a further class-weighted ACDC BDS-Lite set is in Supplementary S5). The first ACDC decoder-only runs had used class weights by mistake, unlike all other retrained ACDC arms; we repeated them without weights and use the repeated runs throughout, which changed no Holm decision. Each ablation and comparator is compared only with retrained references.

The arms follow the central question. \emph{Simple supervision} is U-Net with the boundary decoder and $L_{\mathrm{boundary}}$ only. \emph{The transfer path} (distillation and gate) is added to that arm, giving BDS-Lite without the surface loss, and to the decoder with the surface loss, giving full BDS-Lite; BDS-Lite without distillation ($\lambda_d=0$) and without the gate ($\alpha=0$) separate its parts. \emph{The surface loss} is also tested alone on U-Net. With U-Net these give eight cells, of which U-Net, surface only, decoder only, and decoder plus surface form a two-by-two factorial (Supplementary S3).

The comparators ask whether something simpler or more standard would do as well. U-Net+GSL \citep{Celaya2024A} and U-Net+Boundary Loss \citep{Kervadec2019Boundary} test whether one added loss suffices; Attention U-Net \citep{Oktay2018Attention} at the same width is a standard architectural alternative in 2D; and LB-UNet \citep{Xu2024LBUNet} is the closest lightweight boundary-assisted network. We trained all of them under the protocol above rather than their original ones. LB-UNet needed multi-class output heads, and Boundary Loss uses our region term; the adaptations, loss equations, and source commits are in Supplementary S2--S3.

\subsection*{Additional trainings}
Three seeds left the central contrasts unresolved, so further runs were added under a frozen protocol. The protocol, its analysis code, and its configurations were committed with source hashes before the runs started (\path{revision/EXPANSION_PROTOCOL.md}). It was written after the seeds-1--3 results were known, which is why confirmation uses new seeds only. The recipes were unchanged, only the seed differed, and nothing was tuned. Seed 1 of the four arms rebuilt for this purpose reproduced the archived validation DSC exactly over four epochs before training began.

The protocol added seeds 4--10 of U-Net, decoder-only, BDS-Lite, and U-Net+GSL on ISIC2018, seeds 1--3 of the decoder-plus-surface cell on all datasets, and the corrected ACDC decoder runs (40 runs). All runs passed checks of completeness, hashes, and checkpoint selection, with independent recomputation of test metrics on a sample of runs, before any of their test metrics was read.

\subsection*{Evaluation metrics}
Predictions are the argmax over classes, with no thresholding, test-time augmentation, or post-processing. Per image and foreground class we compute DSC, IoU, Boundary F1 (tolerance 2\,px), and HD95 and ASSD from bidirectional surface distances on the $224\times224$ grid, in pixels. A class absent from both prediction and reference is undefined and excluded. A missed structure or a false alarm scores zero on DSC, IoU, and Boundary F1 and has undefined surface distances, so primary HD95 and ASSD omit such cases. We therefore report missed-structure and false-alarm rates separately, together with the registered all-case sensitivity analysis, which assigns the image diagonal to one-sided empty pairs and zero to both-empty pairs. Values are macro-averaged over classes per image. This convention differs from the original submission, which counted classes absent from both masks as perfect; Supplementary S4 gives the original values.

For ACDC and Synapse we also stack slice predictions into volumes and compute 3D DSC, HD95, and ASSD per patient or case, without voxel spacing, as in the TransUNet evaluator. Millimeter distances and ACDC ejection-fraction and myocardial-mass errors require voxel spacing and end-diastolic and end-systolic frame assignment, which we could not verify for the processed data, so they are not reported. On ISIC2018 we also report Jaccard and thresholded Jaccard at original resolution following the challenge protocol \citep{Codella2019Skin}. Deployment cost comprises parameters, FLOPs, median latency of 200 timed passes after 50 warm-up passes on an idle GPU (batch sizes 1, 8, 16) and on CPU, and CUDA kernel counts from \texttt{torch.profiler}.

\subsection*{Statistical analysis}
We keep two questions apart: whether particular trained models differ on a test set, and whether training recipes differ. The second is what the central question requires, and it is the primary analysis wherever several runs allow it.

\emph{Recipe-level analysis.} The unit of replication is the training run, and equal seed numbers do not pair runs of different arms. For each contrast, $\Delta$ is the difference of arm means of run means, oriented so that positive favors the first arm. Its standard error combines the Welch standard error of run means with the standard error over test units of the per-unit difference, and intervals use the $t$ distribution with Welch degrees of freedom.

The frozen protocol has six contrasts among the four ten-seed ISIC2018 arms, with Boundary F1 as the primary endpoint. Confirmation uses fresh seeds 4--10, Holm-corrected over the six contrasts within each metric \citep{Holm1979A}, because seeds 1--3 prompted the questions; effect sizes use all ten seeds. A contrast is \emph{confirmed} if its fresh-seed Holm $p<0.05$, its interval excludes zero, and its sign matches the all-seed estimate; \emph{supported, not confirmed} if only the all-seed interval excludes zero; and otherwise \emph{not established}.

On all three datasets the protocol also fixed, as secondary analyses with seeds 1--3, the simple effects of the factorial and the contrasts of the transfer path with and without the surface loss. These are estimated per dataset and pooled by inverse variance with Cochran's $Q$. The factorial interaction is $(\text{decoder+surface}-\text{decoder})-(\text{surface}-\text{U-Net})$. These contrasts carry effect labels but no confirmation status. Additive component effects over the eight cells, which separate distillation from the gate, are exploratory.

Smallest effect sizes of interest (SESOI) are engineering thresholds, not clinical ones: 0.01 for DSC, IoU, and Boundary F1, 2\,px for HD95, and 0.5\,px for ASSD. An interval wholly beyond the threshold is a relevant gain or harm; one that excludes zero but overlaps the threshold is a gain or harm of unclear size; one wholly inside $\pm$SESOI is no relevant effect; and one that includes zero and extends past the threshold is undetermined. Only ``no relevant effect'' is evidence of absence.

\emph{Model-level analysis.} For the three-seed arms we report the families registered in the repository (\path{analysis/revision/prereg_revision.md}) before the new results were generated, but after the original results were known. The statistical unit is the image (ISIC2018, $n=1{,}000$), patient (ACDC, $n=20$), or case (Synapse, $n=12$). Per unit, metrics are averaged over seeds and slices, and differences are paired on samples where the metric is defined in both arms. We report mean differences with 95\% percentile bootstrap intervals (10,000 resamples) and two-sided Wilcoxon signed-rank tests with Holm correction within families: BDS-Lite versus U-Net (per model set, five metrics), each ablation versus BDS-Lite, each ablation versus U-Net, and each comparator versus U-Net and versus BDS-Lite. A Friedman test across all retrained arms is in Supplementary S7 \citep{Demsar2006Statistical}. These intervals describe fixed trained models and exclude retraining, so they are descriptive when comparing recipes. The approximate minimum detectable paired effect (80\% power, unadjusted) on the original models is \R{MDE.acdc.dsc} DSC and \R{MDE.acdc.hd95}\,px HD95 on ACDC, and \R{MDE.synapse.dsc} DSC and \R{MDE.synapse.hd95}\,px HD95 on Synapse.

\emph{Post hoc additions.} For three-seed arms we added, after seeing results, a run-level margin: for each dataset and metric we pool the variance of the three run means within arms (all arms except LB-UNet, whose scale differs) and take the 95\% half-width for a difference of two three-run means, $t_{0.975,\mathrm{df}}\,s\sqrt{2/3}$. The repeatability analysis and the measurements of distillation and the gate were also designed after the results were known.

%% =====================================================================
\section*{Results}

\subsection*{How large an effect can be resolved?}
Any answer depends on how far a recipe's result moves when it is simply trained again, so we measured that first. On ISIC2018 the standard deviation of Boundary F1 between runs of one recipe was \R{XM.decoder.boundary_f1.sd} to \R{XM.gsl.boundary_f1.sd} (Table~\ref{tab:tenseed}), larger than the 0.01 threshold and similar to the differences of interest. With ten runs per arm, intervals on Boundary F1 contrasts were still about $\pm0.02$ wide.

Image-level statistics do not capture this variation. We split the ten runs of one recipe into two disjoint groups, averaged each image's metric within each group, and applied the test used in our registered families, a Wilcoxon signed-rank test on the 1,000 per-image differences. There is no true difference between recipes in this comparison. With three runs per group, $p<0.05$ occurred in \R{SS.same.min.3.boundary_f1.p05}--\R{SS.same.max.3.boundary_f1.p05}\% of splits for Boundary F1, depending on the recipe, and in \R{SS.same.min.3.dsc.p05}--\R{SS.same.max.3.dsc.p05}\% for DSC and \R{SS.same.min.3.hd95.p05}--\R{SS.same.max.3.hd95.p05}\% for HD95 (Supplementary S14). Five runs per group did not help (\R{SS.same.min.5.boundary_f1.p05}--\R{SS.same.max.5.boundary_f1.p05}\% for Boundary F1). The difference between groups exceeded the 0.01 threshold in \R{SS.same.mean.3.boundary_f1.beyond}\% of three-run splits on average. The test is not faulty: these particular trained models do differ on these images. That is a different question from whether two recipes differ, and more test images help only with the first. This post hoc analysis draws on ten runs per arm, so its percentages are indicative, but it is why the answers below rest on run-level intervals.

\begin{table}[ht]
\centering
\resizebox{\linewidth}{!}{\input{generated/tab_ten_seed.tex}}
\caption{\label{tab:tenseed}ISIC2018 official test set, ten independent trainings per arm. Mean $\pm$ SD of run means. HD95 and ASSD in pixels. Seeds 1--3 are the retrained runs; seeds 4--10 were added under the frozen protocol.}
\end{table}

\begin{table}[ht]
\centering
\resizebox{\linewidth}{!}{\input{generated/tab_ten_seed_contrasts.tex}}
\caption{\label{tab:tencontrasts}Frozen contrasts on ISIC2018 Boundary F1, the primary endpoint. Positive favors the first arm. The interval combines variation between runs and between test images (Statistical analysis). Holm correction is over the six contrasts. Status rules were fixed before the runs. All five metrics are in Supplementary S13.}
\end{table}

\begin{figure}[ht]
\centering
\includegraphics[width=\linewidth]{../figures/ten_seed.pdf}
\caption{\label{fig:tenseed}Ten independent trainings per arm on the ISIC2018 official test set. Open markers: seeds 1--3, which prompted the questions. Filled markers: seeds 4--10, used for confirmation. Vertical bars: means over ten runs. The HD95 axis is reversed, so better values lie to the right in every panel.}
\end{figure}

\subsection*{Does training-only boundary supervision help?}
U-Net with the auxiliary decoder had the highest ISIC2018 Boundary F1 of the four ten-seed arms (\R{XM.decoder.boundary_f1.mean}, against \R{XM.unet.boundary_f1.mean} for U-Net; Fig.~\ref{fig:tenseed}). The difference was \R{X.decoder.unet.boundary_f1.all.delta} \R{X.decoder.unet.boundary_f1.all.ci} (Table~\ref{tab:tencontrasts}). By the rule fixed in advance this is supported but not confirmed, because the fresh-seed estimate, although of similar size (\R{X.decoder.unet.boundary_f1.fresh.delta}, \R{X.decoder.unet.boundary_f1.fresh.ci}), did not survive Holm correction ($p=$\,\R{X.decoder.unet.boundary_f1.fresh.holm}). DSC, IoU, HD95, and ASSD showed no relevant effect, as expected if the change is confined to contours (Eq.~\ref{eq:firstorder}).

The change looks like better contour placement rather than more contour. Boundary precision rose from \R{BP.M0.precision} to \R{BP.M2.precision} and recall from \R{BP.M0.recall} to \R{BP.M2.recall}, and spurious components did not increase (\R{BP.M0.spurious} and \R{BP.M2.spurious} per image). Other datasets point the same way with less certainty. In the seeds-1--3 factorial the decoder changed Boundary F1 by \R{I.decoder-without-surface.isic2018.boundary_f1.three.est}, \R{I.decoder-without-surface.acdc.boundary_f1.three.est}, and \R{I.decoder-without-surface.synapse.boundary_f1.three.est} on ISIC2018, ACDC, and Synapse, with every individual interval including zero and a pooled estimate of \R{I.decoder-without-surface.pooled.boundary_f1.three.est} \R{I.decoder-without-surface.pooled.boundary_f1.three.ci}, a gain of unclear size (Table~\ref{tab:central}, Fig.~\ref{fig:central}). The ACDC value is measured against a retrained U-Net that is a low reference (see below), so it should be read with caution. Pooled DSC and HD95 showed no relevant effect. This arm deploys as a plain U-Net.

The defensible statement is that training-only boundary supervision had the highest ISIC2018 Boundary F1 of the ten-run arms, pointed the same way on every dataset, and costs nothing at inference, but that its gain over U-Net was not confirmed on seeds that played no part in suggesting it.

\subsection*{Does distillation add to the decoder?}
This is the central contrast, and on Boundary F1 it gave no established benefit. Adding distillation and the gate to the decoder changed Boundary F1 by \R{I.distillationandgate-without-surface.isic2018.boundary_f1.three.est}, \R{I.distillationandgate-without-surface.acdc.boundary_f1.three.est}, and \R{I.distillationandgate-without-surface.synapse.boundary_f1.three.est} on ISIC2018, ACDC, and Synapse (pooled \R{I.distillationandgate-without-surface.pooled.boundary_f1.three.est} \R{I.distillationandgate-without-surface.pooled.boundary_f1.three.ci}). Adding them to the decoder with the surface loss changed it by \R{I.distillationandgate-with-surface.isic2018.boundary_f1.three.est}, \R{I.distillationandgate-with-surface.acdc.boundary_f1.three.est}, and \R{I.distillationandgate-with-surface.synapse.boundary_f1.three.est} (pooled \R{I.distillationandgate-with-surface.pooled.boundary_f1.three.est} \R{I.distillationandgate-with-surface.pooled.boundary_f1.three.ci}). Both pooled intervals include zero and extend past the threshold, so both are undetermined (Table~\ref{tab:central}, Fig.~\ref{fig:central}). They do bound the answer: pooled gains above about 0.02 are unlikely. The first estimate leans positive and a small gain of about 0.01 cannot be excluded, but ISIC2018 does not share that direction, and the per-dataset intervals are wide. For DSC and HD95 the pooled label was no relevant effect in both contrasts, which is evidence of absence at those thresholds.

The additive model, which is exploratory but uses all eight cells, separates the two parts (Table~\ref{tab:components}). Pooled Boundary F1 effects were \R{K.pooled.boundary_f1.distill.est} \R{K.pooled.boundary_f1.distill.ci} for distillation and \R{K.pooled.boundary_f1.gate.est} \R{K.pooled.boundary_f1.gate.ci} for the gate, again undetermined but close to zero, with no relevant effect on DSC or HD95. The registered image-level family agrees in direction: none of the \R{LOO.tests} leave-one-out contrasts showed the full recipe to be better, and \R{LOO.better} favored a reduced recipe, all on ISIC2018. Varying $\alpha$ over $\{0,0.1,0.25,0.5,1\}$ on ISIC2018 moved no endpoint beyond the threshold relative to $\alpha=0.25$ (\R{SENS.isic2018.alpha.M6.equivalent} of \R{SENS.isic2018.alpha.M6.tests} contrasts within it, the others inconclusive), and a single-seed check of $\lambda_b=\lambda_d\in\{0.01,0.2\}$ gave no support for $0.05$ as a favorable value (Supplementary S11).

The arm behind the first estimate deserves a direct comment. BDS-Lite without the surface loss had the highest three-seed Boundary F1 of the ablation arms on ACDC and Synapse, and its ACDC difference from U-Net passed Holm correction and exceeded the run-level margin (Table~\ref{tab:arms}). We do not treat this as a positive finding. Its margin over the decoder, the reference that matters here, is what Table~\ref{tab:central} shows; the ACDC U-Net is a low reference; its Synapse difference depends on how samples are weighted (Supplementary S5); and choosing a variant from test results would be circular. It is a candidate for testing on new seeds, not a result.

\begin{table}[ht]
\centering
\resizebox{\linewidth}{!}{\input{generated/tab_central.tex}}
\caption{\label{tab:central}The central contrasts on all datasets, seeds 1--3, with 95\% intervals. Transfer: distillation and gate added to the decoder, without the surface loss (BDS-Lite without surface loss versus decoder) or with it (BDS-Lite versus decoder plus surface loss). Positive favors the arm with more components (HD95 sign reversed, in pixels). Pooled by inverse variance. Labels are the frozen ones; an undetermined interval includes zero and extends past the threshold. On ISIC2018 the decoder-versus-U-Net contrast also has the ten-seed estimate in Table~\ref{tab:tencontrasts}.}
\end{table}

\begin{figure}[ht]
\centering
\includegraphics[width=\linewidth]{../figures/central_question.pdf}
\caption{\label{fig:central}The central question on Boundary F1. From top: training-only supervision against U-Net, the transfer path added to the decoder without and with the surface loss, and the full recipe against the decoder alone. Lines are 95\% intervals that include retraining variation. The gray band is the $\pm0.01$ threshold. Pooled rows here are dominated by ACDC, which has the narrowest intervals.}
\end{figure}

\begin{table}[ht]
\centering
\resizebox{\linewidth}{!}{\input{generated/tab_components.tex}}
\caption{\label{tab:components}Exploratory additive component effects with 95\% intervals, fitted over the eight ablation cells with the ISIC2018 ten-seed arms. Positive favors adding the component (HD95 sign reversed, in pixels). Labels are the frozen ones. The additive form is an assumption, since the interaction is unresolved.}
\end{table}

\begin{table}[ht]
\centering
\resizebox{\linewidth}{!}{\input{generated/tab_all_arms.tex}}
\caption{\label{tab:arms}All retrained arms relative to the retrained U-Net, seeds 1--3: oriented mean differences (positive favors the arm; HD95 in pixels with the sign reversed). ISIC2018: official test set. $^{*}$Holm $p<0.05$ within family and dataset with an interval excluding zero, which concerns fixed models and not recipes (see ``How large an effect can be resolved?''). $^{\dagger}$The difference between run means exceeds the run-level margin. The retrained ACDC U-Net is a low reference (see text). Further contrasts and the individual runs are in Supplementary S5--S7.}
\end{table}

\subsection*{Why the transfer added so little}
Measurements of the trained networks suggest why. They are post hoc and use inference only (Supplementary S13).

Distillation worked as a transfer mechanism. Applying the frozen teacher head to the deployed projection $q_s$ gave a boundary average precision of \R{MECH.isic2018.head.dist}, \R{MECH.acdc.head.dist}, and \R{MECH.synapse.head.dist} in distilled arms against \R{MECH.isic2018.head.plain}, \R{MECH.acdc.head.plain}, and \R{MECH.synapse.head.plain} in undistilled arms on ISIC2018, ACDC, and Synapse (Welch $p\le$\,\R{MECH.acdc.head.p}), although well below the teacher itself (\R{MECH.isic2018.teacher}, \R{MECH.acdc.teacher}, \R{MECH.synapse.teacher}).

The signal did not detectably reach the features that produce the prediction, which is the doubt raised in Methods. A linear readout of boundaries from $F_s$ differed between distilled and undistilled arms by \R{MECH.isic2018.fs.diff} \R{MECH.isic2018.fs.ci}, \R{MECH.acdc.fs.diff} \R{MECH.acdc.fs.ci}, and \R{MECH.synapse.fs.diff} \R{MECH.synapse.fs.ci}, and arms without a decoder already supported a readout of \R{MECH.isic2018.fs.none}, \R{MECH.acdc.fs.none}, and \R{MECH.synapse.fs.none}. Boundary information was thus already readable from $F_s$ without a boundary decoder, so a $1\times1$ projection can satisfy the loss by reading it. Consistent with this, the distillation gradient was nearly orthogonal to the segmentation gradient (mean cosine \R{GR.isic2018.distill.all.cos}, \R{GR.acdc.distill.all.cos}, and \R{GR.synapse.distill.all.cos}) and at most \R{GR.isic2018.distill.all.ratio} of its norm in seed-1 probes.

The gate offered no second route. It ran at \R{GATE.isic2018.bound}\%, \R{GATE.acdc.bound}\%, and \R{GATE.synapse.bound}\% of its bound, so it was not close to the identity, but most of that was a fixed per-channel scale (\R{GATE.isic2018.const} of \R{GATE.isic2018.abs}, \R{GATE.acdc.const} of \R{GATE.acdc.abs}, and \R{GATE.synapse.const} of \R{GATE.synapse.abs}), which the output convolution can absorb. It was not selective for the boundary band (ratio of band to off-band modulation \R{MECH.isic2018.gate.band}, \R{MECH.acdc.gate.band}, \R{MECH.synapse.gate.band}). Removing it from trained models changed fewer than \R{GATE.labels.max}\% of pixel labels, and setting it to the identity on the original models changed DSC, IoU, and Boundary F1 by at most \R{G.max.overlap}.

\subsection*{The full recipe and its surface loss}
\emph{Against the decoder alone.} Adding the surface loss, distillation, and the gate together lowered ISIC2018 Boundary F1. This is the clearest contrast in the study: \R{X.full.decoder.boundary_f1.all.delta} \R{X.full.decoder.boundary_f1.all.ci} over ten seeds and \R{X.full.decoder.boundary_f1.fresh.delta} on fresh seeds (Holm $p=$\,\R{X.full.decoder.boundary_f1.fresh.holm}), confirmed, with the interval beyond the threshold. Boundary precision and recall both fell relative to the decoder alone (\R{BP.M6.precision} and \R{BP.M6.recall}). On the other four metrics the intervals included zero.

\emph{Attribution.} This confirmed deficit cannot be assigned to one part. The step from the decoder to BDS-Lite has two stages, and on ISIC2018 neither is resolved alone: adding the surface loss to the decoder gave \R{I.surface-with-decoder.isic2018.boundary_f1.ten.est} \R{I.surface-with-decoder.isic2018.boundary_f1.ten.ci}, and then adding the transfer path gave \R{I.distillationandgate-with-surface.isic2018.boundary_f1.three.est} \R{I.distillationandgate-with-surface.isic2018.boundary_f1.three.ci}. The additive model assigns most of it to the surface loss (ISIC2018 \R{K.isic2018.boundary_f1.surface.est} \R{K.isic2018.boundary_f1.surface.ci}; distillation \R{K.isic2018.boundary_f1.distill.est} \R{K.isic2018.boundary_f1.distill.ci}), but that rests on additivity, and the interaction was unresolved on every dataset and pooled (\R{I.interaction.pooled.boundary_f1.three.est} \R{I.interaction.pooled.boundary_f1.three.ci}; ISIC2018 with ten-seed U-Net and decoder cells, \R{I.interaction.isic2018.boundary_f1.ten.est} \R{I.interaction.isic2018.boundary_f1.ten.ci}), as we had anticipated with three seeds per cell (Supplementary S13).

\emph{Surface loss.} It did not reduce HD95, the purpose for which it was included: the additive model gave no relevant effect on any dataset or pooled (\R{K.pooled.hd95.surface.est}\,px \R{K.pooled.hd95.surface.ci}). Evidence that it lowers Boundary F1 is concentrated on ISIC2018. Added to U-Net there, it changed Boundary F1 by \R{I.surface-without-decoder.isic2018.boundary_f1.ten.est} \R{I.surface-without-decoder.isic2018.boundary_f1.ten.ci} and increased distant false positives from \R{EB.isic2018.M0.fp_far} to \R{EB.isic2018.M1.fp_far} pixels per image. The pooled additive effect was \R{K.pooled.boundary_f1.surface.est} \R{K.pooled.boundary_f1.surface.ci}, but the direct seeds-1--3 factorial cells did not show it (pooled \R{I.surface-without-decoder.pooled.boundary_f1.three.est} \R{I.surface-without-decoder.pooled.boundary_f1.three.ci} without and \R{I.surface-with-decoder.pooled.boundary_f1.three.est} \R{I.surface-with-decoder.pooled.boundary_f1.three.ci} with the decoder). In seed-1 probes its gradient norm relative to the segmentation gradient had a median of \R{GR.isic2018.sdf.all.ratio} on ISIC2018 but \R{GR.acdc.sdf.all.ratio} and \R{GR.synapse.sdf.all.ratio} on ACDC and Synapse, so it may have had little influence on the multi-structure tasks (Supplementary S10); the normalization in Eq.~\ref{eq:sdf} is a plausible but untested cause.

\emph{Against U-Net.} Ten independent trainings per arm gave no established Boundary F1 difference between BDS-Lite and U-Net on ISIC2018: \R{X.full.unet.boundary_f1.all.delta} \R{X.full.unet.boundary_f1.all.ci}, with fresh seeds alone giving \R{X.full.unet.boundary_f1.fresh.delta}. DSC and HD95 showed no relevant effect (\R{X.full.unet.dsc.all.delta} \R{X.full.unet.dsc.all.ci}; \R{X.full.unet.hd95.all.delta}\,px \R{X.full.unet.hd95.all.ci}). The HD95 benefit claimed in the original submission therefore does not hold on independent data.

Table~\ref{tab:twosets} sets this beside the earlier model sets, which illustrate why run-level inference matters. The original models differed in Boundary F1 by \R{P.isic2018.S6.S0.boundary_f1.delta} (\R{P.isic2018.S6.S0.boundary_f1.ci}; Holm $p=$\,\R{P.isic2018.S6.S0.boundary_f1.holm}), and the difference persisted after excluding all \R{DUP.n.0.9} test images with a possible repeat in the development pool (\R{DUP.P.boundary_f1.excluding_r_gt_0.9.delta}, \R{DUP.P.boundary_f1.excluding_r_gt_0.9.ci}). Three retrained runs per arm gave \R{P-matched.isic2018.M6.M0.boundary_f1.delta}, and ten gave the value above. The original value lies above the central 95\% of differences between three retrained runs per arm (\R{SS.contrast.full.unet.3.boundary_f1.range}), so sampling of runs is an unlikely full explanation; those models also used another partition and a trainer that no longer exists. Their interval, narrow as it is, describes three fixed models and not the recipe.

On ACDC, HD95 and ASSD moved against BDS-Lite in the original models (\R{P.acdc.S6.S0.hd95.delta}\,px, \R{P.acdc.S6.S0.hd95.ci}; \R{P.acdc.S6.S0.assd.delta}\,px, \R{P.acdc.S6.S0.assd.ci}) without surviving Holm correction. In the retrained set BDS-Lite was higher in DSC (\R{P-matched.acdc.M6.M0.dsc.delta}, \R{P-matched.acdc.M6.M0.dsc.ci}), which reflects a low retrained U-Net at least as much as any gain. The two BDS-Lite sets were close (same-recipe gap \R{RL.acdc.dsc.bgap}), whereas the two U-Net sets differed by \R{RL.acdc.dsc.ugap}, more than the margin of $\pm$\R{RL.acdc.dsc.margin}, for reasons that cannot be established because the original trainer is lost (Supplementary S5). The transfer-path contrasts do not involve this U-Net.

On Synapse every slice-level interval included zero in both sets (DSC \R{P.synapse.S6.S0.dsc.delta}, \R{P.synapse.S6.S0.dsc.ci}; \R{P-matched.synapse.M6.M0.dsc.delta}, \R{P-matched.synapse.M6.M0.dsc.ci}). At volume level the two sets again differed: 3D DSC changed by \R{V.P.synapse.dsc.delta} \R{V.P.synapse.dsc.ci} for the original models and \R{V.P-replication.synapse.dsc.delta} (\R{V.P-replication.synapse.dsc.ci}; Holm $p=$\,\R{V.P-replication.synapse.dsc.holm}) for the retrained models, and no 3D HD95 or ASSD difference was significant (Supplementary S9). The overlap--distance conflict on Synapse highlighted in the original submission was therefore not established.

\begin{table}[ht]
\centering
\resizebox{\linewidth}{!}{\input{generated/tab_two_sets.tex}}
\caption{\label{tab:twosets}BDS-Lite versus U-Net in the two three-seed model sets. Mean $\pm$ SD over three seeds of image- or slice-level values; $\Delta$ is the unit-weighted paired difference with distance differences negated, so positive favors BDS-Lite, and it need not equal the difference of the means shown. ISIC2018: official test set; the ten-seed results in Table~\ref{tab:tencontrasts} supersede the retrained row. $^{*}$Holm $p<0.05$ with an interval excluding zero. Run-level margin: see Statistical analysis. Values differ from the original Table~2 because classes absent from both masks are no longer scored as perfect.}
\end{table}

\subsection*{Other boundary methods and published results}
U-Net+GSL shows how three seeds can mislead, and it gives the decoder its second confirmed comparison. With seeds 1--3 GSL was the most consistent comparator: relative to U-Net, Boundary F1 changed by \R{C.isic2018.GSL.M0.boundary_f1.delta}, \R{C.acdc.GSL.M0.boundary_f1.delta}, and \R{C.synapse.GSL.M0.boundary_f1.delta} on ISIC2018, ACDC, and Synapse, all \R{GSLM6.cells} contrasts with BDS-Lite favored it, and \R{GSLM6.holm} survived Holm correction. With ten seeds on ISIC2018 that lead vanished: GSL differed from BDS-Lite by \R{X.gsl.full.boundary_f1.all.delta} \R{X.gsl.full.boundary_f1.all.ci} and from U-Net by \R{X.gsl.unet.boundary_f1.all.delta} \R{X.gsl.unet.boundary_f1.all.ci}, and its own run SD was the largest of the four arms (\R{XM.gsl.boundary_f1.sd}). The decoder alone exceeded it, confirmed (\R{X.gsl.decoder.boundary_f1.all.delta} \R{X.gsl.decoder.boundary_f1.all.ci}; fresh Holm $p=$\,\R{X.gsl.decoder.boundary_f1.fresh.holm}). On DSC and HD95 GSL showed no relevant effect against U-Net. These conclusions concern ISIC2018; on ACDC and Synapse GSL has three seeds (Table~\ref{tab:arms}).

Resampling the ten-seed arms shows how common such reversals are. Of random draws of three runs per arm, \R{SS.contrast.full.unet.3.boundary_f1.same}\% gave the same sign as the ten-seed BDS-Lite versus U-Net estimate, and \R{SS.contrast.full.unet.3.boundary_f1.opp}\% gave $p<0.05$ in the opposite direction; for GSL versus BDS-Lite, whose ten-seed estimate is near zero, \R{SS.contrast.gsl.full.3.boundary_f1.p05}\% of draws reached $p<0.05$. The large difference was stable, however: BDS-Lite versus decoder kept its sign in \R{SS.contrast.full.decoder.3.boundary_f1.same}\% of draws. Results that passed Holm correction at three seeds, including the original Boundary F1 difference and the lead of GSL, thus did not survive additional training, whereas BDS-Lite versus decoder held its direction in every draw.

U-Net+Boundary Loss stayed close to U-Net on ISIC2018 except for HD95 (\R{C.isic2018.C1.M0.hd95.delta}\,px, \R{C.isic2018.C1.M0.hd95.ci}), and its lower ACDC HD95 was the only contrast with BDS-Lite, apart from LB-UNet, to exceed the run-level margin (\R{C-vs-M6.acdc.C1.M6.hd95.delta}\,px, \R{C-vs-M6.acdc.C1.M6.hd95.ci}). Attention U-Net stayed close to U-Net in overlap.

LB-UNet had the largest effects in both directions. On ISIC2018 it changed DSC by \R{C.isic2018.C4.M0.dsc.delta} \R{C.isic2018.C4.M0.dsc.ci} and HD95 by \R{C.isic2018.C4.M0.hd95.delta}\,px \R{C.isic2018.C4.M0.hd95.ci}, beyond the run-level margin, while changing Boundary F1 by \R{C.isic2018.C4.M0.boundary_f1.delta}. This is the clearest metric-family disagreement we observed. On the multi-class tasks its DSC changed by \R{C.acdc.C4.M0.dsc.delta} and \R{C.synapse.C4.M0.dsc.delta}; training converged, so this reflects the adapted network under this protocol, whereas LB-UNet was designed for 256-pixel binary lesions with augmentation. A Friedman test across all arms rejected equal performance for every dataset and metric (Supplementary S7).

Table~\ref{tab:lit} places the compact setting next to published results obtained under other protocols. Absolute accuracy lies below heavier systems, as expected for a 1.9\,M-parameter network trained without augmentation or pretraining. These numbers are context, not a ranking.

\begin{table}[ht]
\centering
\resizebox{0.92\linewidth}{!}{\input{generated/tab_lit_compact.tex}}
\caption{\label{tab:lit}Published results and this study. Not directly comparable: protocols differ in training data, pretraining (ImageNet for TransUNet and Swin-Unet), augmentation, and resolution. Published values were verified against the cited tables. Synapse and ACDC rows are volume-level; HD95 is in index units for this study and computed without spacing in the TransUNet code. ISIC2018 rows use the same official test set, but challenge entries had no restriction on models or data. A longer table is in Supplementary S7.}
\end{table}

\subsection*{Which errors change}
Most error pixels on ISIC2018 were not near the contour, which helps explain why a boundary signal moves some metrics and not others. Only \R{ERR.isic2018.S0.near}--\R{ERR.isic2018.M0.near}\% of wrong U-Net pixels lay within 2\,px of the reference boundary, compared with \R{ERR.acdc.S0.near}\% on ACDC and \R{ERR.synapse.M0.near}--\R{ERR.synapse.S0.near}\% on Synapse. Most ISIC2018 errors are lesion area missed or added in bulk, so a band-limited signal can change contour agreement while leaving most of the error untouched.

Spurious components without reference overlap were higher under BDS-Lite than U-Net in all six dataset and three-seed model-set pairs (Supplementary S8), but ten seeds on ISIC2018 weaken this: the means were close (\R{BP.M0.spurious} and \R{BP.M6.spurious}), so the pattern should not be read as established there. On ACDC false alarms on slices without the structure rose (\R{CB.acdc.S0.empty-reference.falsealarm}\% to \R{CB.acdc.S6.empty-reference.falsealarm}\%; \R{CB.acdc.M0.empty-reference.falsealarm}\% to \R{CB.acdc.M6.empty-reference.falsealarm}\%), and distant false negatives fell (\R{EB.acdc.S0.fn_far} to \R{EB.acdc.S6.fn_far}; \R{EB.acdc.M0.fn_far} to \R{EB.acdc.M6.fn_far} pixels per slice), a shift toward predicting more foreground rather than better contours. LB-UNet produced \R{EB.isic2018.C4.spurious} spurious components per ISIC2018 image, consistent with its lower HD95, since one remote component can set the distance tail.

Small structures were not protected by per-class Dice. On Synapse, reference objects below 1\% of the image were missed entirely in \R{CB.synapse.S0.lt1pct.missed}\% and \R{CB.synapse.M0.lt1pct.missed}\% of instances by the two U-Nets and \R{CB.synapse.S6.lt1pct.missed}\% and \R{CB.synapse.M6.lt1pct.missed}\% by the two BDS-Lite sets, with none missed above 5\% of the image.

Such cases also affect the distance metrics. Primary HD95 and ASSD skip missed structures and false alarms, so a model can look better by missing more. In the all-case analysis, the retrained ACDC HD95 difference changed sign, from \R{P-matched.acdc.M6.M0.hd95.delta}\,px to \R{AC.acdc.retrained.hd95.delta}\,px \R{AC.acdc.retrained.hd95.ci}, and the Synapse differences also turned positive in both sets, with intervals including zero (Supplementary S8). Whether BDS-Lite worsens ACDC distances, which the original submission emphasized, thus depends on how undefined cases are scored.

Fig.~\ref{fig:qualitative} shows ISIC2018 examples selected by a fixed rule. The largest per-image HD95 differences came from interior holes and remote components rather than contour placement. In the largest-gain case, U-Net and U-Net+GSL left holes inside the lesion (HD95 \R{QUAL.isic2018.gain.S0.hd95} and \R{QUAL.isic2018.gain.GSL.hd95}\,px, against \R{QUAL.isic2018.gain.S6.hd95}\,px for BDS-Lite); in the largest-loss case BDS-Lite predicted \R{QUAL.isic2018.loss.S6.spurious} remote components and LB-UNet none. In the median case all models differed by a thin rim. Extremes were unstable: in \R{QUAL.extreme.unstable} of the \R{QUAL.extreme.n} extreme cases across datasets, at least one seed showed a difference below 2\,px in the selected direction or the opposite sign (ACDC and Synapse examples: Supplementary S12).

\begin{figure}[ht]
\centering
\includegraphics[width=\linewidth]{../figures/qualitative_isic2018.pdf}
\caption{\label{fig:qualitative}Official ISIC2018 test images with the largest positive, median, and largest negative seed-averaged HD95 difference between the original U-Net and BDS-Lite models (positive: BDS-Lite better); the row label gives the difference and the displayed seed, whose difference is closest to the seed mean. Upper rows: reference contour (green) and prediction (magenta). Lower rows: false positives (red) and false negatives (blue). All other columns are retrained models. Images: ISIC archive, CC-0. Identifiers and per-panel statistics: \texttt{revision/figures/qualitative\_manifest\_v2.json}.}
\end{figure}

\subsection*{Deployment cost}
The transfer path has a measured price. Deployed BDS-Lite has 1,072 more parameters than U-Net (0.056\%) and 0.68\% more FLOPs (Table~\ref{tab:resource}), whereas the decoder-only and GSL arms deploy as plain U-Nets. A small parameter count says nothing about effectiveness, and the overhead is not free in time: latency rose by \R{RESX.BDSLitedeployed.gpu_fp32_b1_ms_median.pct}\% at batch size 1 and \R{RESX.BDSLitedeployed.gpu_fp32_b16_ms_median.pct}\% at batch size 16 on the GPU, and \R{RESX.BDSLitedeployed.cpu_ms_b1_1thread.pct}\% on one CPU thread. The gate adds CUDA kernels (\R{RESX.UNet.kernels} to \R{RESX.BDSLitedeployed.kernels}) and full-resolution element-wise work, which FLOP counts do not capture. The 24\% batch-1 overhead of the original submission was not reproduced. Replacing the gate by the identity restores the U-Net graph exactly (\R{RESX.BDSLitegatecollapsed.kernels} kernels). Training is costlier: in the original profiling, one step took \R{TRAIN.bds.ms}\,ms and \R{TRAIN.bds.mb}\,MB for BDS-Lite, \R{TRAIN.unet.ms}\,ms and \R{TRAIN.unet.mb}\,MB for U-Net, and \R{TRAIN.gsl.ms}\,ms and \R{TRAIN.gsl.mb}\,MB for U-Net+GSL (batch size 8). LB-UNet, with about 1\% of U-Net's FLOPs, was slower on the GPU at batch size 1 but far faster on one CPU thread.

\begin{table}[ht]
\centering
\resizebox{\linewidth}{!}{\input{generated/tab_resource.tex}}
\caption{\label{tab:resource}Deployment profile (ISIC2018 input, 2 classes). Latencies are medians in ms (fp32). $^a$Forward-hook count used in the original submission.}
\end{table}

%% =====================================================================
\section*{Discussion}

\subsection*{Does boundary distillation add to training-only supervision?}
In the studied setting we found no established added value. The best-supported boundary result belonged to the decoder alone, which deploys as a plain U-Net. Adding distillation and the gate gave pooled Boundary F1 changes that include zero and no relevant effect on DSC or HD95, and the full recipe scored below the decoder alone on ISIC2018. As for the title's question, we found no condition under which the transfer path clearly helped, and the conditions proposed in our original submission did not hold up.

We do not claim that distillation has no effect or that it is harmful. The pooled intervals leave room for a small Boundary F1 gain near the 0.01 threshold and make one much above 0.02 unlikely, and the confirmed deficit of BDS-Lite belongs to the combination of surface loss, distillation, and gate. Distillation is thus unproven, not refuted, and the conclusion concerns one transfer design in one compact network.

The decoder-only reference is what makes this answer possible. Judged against U-Net, as distillation methods usually are, a recipe containing the decoder inherits whatever the decoder provides, and the transfer path can take the credit. Our original submission made only that comparison.

\subsection*{Why a working transfer left the prediction unchanged}
The measurements show how a mechanism can work and leave accuracy unchanged. Distillation moved boundary-readable signal into a projection, but the features that produce the output held no more of it than those of undistilled arms, and networks without a decoder already held a comparable amount. This matches the concern set out in Methods: a learnable map between student and teacher may absorb the loss, and the premise that $F_s$ lacked boundary information was not borne out at the level of a linear readout. The gate tells a similar story. We intended it to let boundary features act on the output, and it learned a near-constant channel scale instead.

This suggests a design lesson that we did not test: a transfer loss should act on the features the output head reads, without an intermediate map that can satisfy it alone. It also suggests a check that any feature distillation method can report, namely whether the target signal has increased in the features that make the prediction, and not only at the point where the loss is applied. These analyses are post hoc, the readouts are linear, and nonlinear information would be invisible to them, so they explain the accuracy results without proving the explanation.

\subsection*{Why BDS-Lite did not exceed simpler recipes}
The shortfalls of the recipe have identifiable sources. Most ISIC2018 error lies away from the contour, where its added parts do not act. Its transfer path did not change the features that matter. Its surface loss did not reduce the distance tail for which it was included. Its weights were fixed in advance and identical across datasets; tuning might change the picture, although a fair test would need the same search for every arm. This argues for adding components one at a time, each against the arm without it.

Task structure did not organize the results as we first proposed. Our original conclusion, that boundary distillation helps on a single-contour task and induces metric conflict on complex anatomy, rested on the invalid ISIC2018 evaluation and on a Synapse difference that was not established, and three datasets that also differ in modality and sample size could not establish such a dependence in any case. One point of the original analysis does stand: metric families respond to different errors. Boundary F1 responds to agreement near the reference contour and HD95 to whatever lies farthest from it. The decoder's estimated change was confined to the first, and LB-UNet moved them in opposite directions. A report based on either alone would have misjudged these models.

\subsection*{Reliability as part of the answer}
Each evaluation finding determined what could be concluded about distillation. First, a data-partition error produced an apparent HD95 benefit that did not hold on independent images; recording training-set identity in each checkpoint and checking it against the evaluation manifest would have caught it. Second, image-level significance does not answer the question usually asked of it. Identical recipes ``differed'' in most splits on a 1,000-image test set, and a larger test set should make that more frequent. Comparing training methods needs an interval that includes retraining, and our data suggest that three runs per arm are too few for effects near 0.01 to 0.02 in Boundary F1. Third, questions raised by a first set of runs need runs that did not raise them: the confirmed differences from the decoder and the disappearance of the GSL lead both came from that step. Fourth, averaging HD95 over defined cases rewards models that omit structures, so distances should be read with missed-structure and false-alarm rates.

\subsection*{Practical implications}
For a compact segmentation network the results suggest an order of work, worded here by strength of evidence. An auxiliary decoder that is discarded after training is a reasonable first step when boundary agreement matters: it costs nothing at inference and had the best-supported favorable effect, although that effect is not confirmed. Any further boundary component should be judged against that arm rather than against a network with no boundary signal. A surface term should be checked on the target task, because in our experiments it did not reduce HD95. Small effects should be reported with run-level intervals, and a difference that appears at three seeds should be treated as a question and not as a result.

\subsection*{Clinical relevance}
The error types above relate differently to clinical use, although this study makes no clinical-impact claim. In dermoscopy the outline enters border assessment and excision planning \citep{Nachbar1994ABCD}. Because overlap is weakly sensitive to contour shifts there, Boundary F1 is the endpoint closest to this use, and it is the endpoint on which the components differed. In cardiac MRI, ejection fraction and myocardial mass derive from delineated volumes \citep{Bernard2018Deep}. The thin myocardium is highly sensitive to contour shifts (Table~\ref{tab:datasets}), and false-positive components would enter such volumes directly, so the increase in false alarms on ACDC deserves attention. In abdominal CT, contours define organs at risk \citep{Jabbour2014Upper}; small organs were the structures most often missed, and boundary supervision did not change that. Whether any difference reported here is clinically meaningful depends on physical spacing and derived measurements that we could not compute.

\subsection*{Limitations}
The central transfer-path contrasts rest on three runs per arm. Their per-dataset intervals are wide, their pooled Boundary F1 estimates are dominated by ACDC, and they carry effect labels but no fresh-seed confirmation. The decoder gain over U-Net is not confirmed on ISIC2018 and is estimated from three runs per arm on ACDC and Synapse. Most ablation and comparator arms have three seeds in this study. The additional trainings were designed after earlier results were known, and although confirmation used new seeds, the choice of questions was not blind. Test samples are small for ACDC and Synapse. The interaction is unresolved, because three seeds per cell cannot estimate it, so the additive component effects rest on an assumption. The repeatability analysis draws from ten runs per arm and uses one dataset's test set.

One compact backbone and one distillation design were studied, without augmentation, which may inflate remote false positives and may not reflect behavior under standard training. Other transfer designs, teachers, and loss weights might behave differently; $\lambda$ was checked with a single seed. We trained comparators under our protocol rather than their own, which may disadvantage them, particularly LB-UNet. We did not run nnU-Net, whose self-configuring pipeline changes the augmentation, resolution, and schedule held fixed here, nor Swin UNETR, which is natively 3D, nor 2.5D or 3D variants; models process slices independently and nothing enforces inter-slice continuity, which only the volume-level metrics assess. Distances are in pixels, and clinical indices were not computed. The original models cannot be retrained exactly, and their ISIC2018 partition differs from that of all other arms. The Synapse class names used in the supplement are inferred from numeric labels.

%% =====================================================================
\section*{Conclusions}
We asked whether boundary distillation adds to simple training-only boundary supervision in a compact U-Net, and built BDS-Lite so that each of its signals could be removed in turn. In our experiments it showed no established added value. A decoder discarded after training gave the best-supported boundary result, a Boundary F1 gain on ISIC2018 that is supported but not confirmed, at no inference cost. Distillation and the gate gave no established further gain and no relevant pooled change in DSC or HD95, and the full recipe scored below the decoder alone on ISIC2018. Distillation moved boundary signal into a deployed projection without detectably changing the features that make the prediction, which offers an explanation and a testable design lesson. These conclusions depended on comparing recipes over independent runs and on corrected test data: identical recipes usually appeared significantly different on a large test set, and results that had passed multiple-testing correction at three seeds did not survive further training. Boundary methods should be judged against their own training-only supervision, over independent runs, with first-round questions confirmed on new seeds.


\section*{Acknowledgments}
\input{acknowledgments_revision.tex}

\section*{Data and code availability}
Code, configurations, split manifests, per-case results, and analysis scripts are available at \url{https://github.com/Jingchen-Chen/bds-lite}. ISIC2018 Task~1 training, validation and test data are available from the ISIC archive under CC-0; ACDC and Synapse are available from their providers under their terms.

\bibliography{references_zotero}

\end{document}
